%=====================================
%   homological-subdifferential-slide.tex
%   Homological Subdifferential について
%   toshi2019
%=====================================
\documentclass[dvipdfmx,12pt,aspectratio=169,leqno]{beamer}% dvipdfmxしたい
\title{Homological Subdifferential
\subtitle{中間発表}}
\date{2024年11月6日}
\author{大柴 寿浩}





\usetheme{metropolis}
%\usetheme{metroplis}
%\usetheme{Singapore}
\usecolortheme{rose}
\usefonttheme{professionalfonts}
\setbeamertemplate{navigaton symbols}{\false}
\setbeamertemplate{footline}[frame number]
%\usepackage{graphicx,xcolor}





\usepackage{bxdpx-beamer}% dvipdfmxなので必要
\usepackage{pxjahyper}% 日本語で'しおり'したい
\usepackage{minijs}% min10ヤダ
\renewcommand{\kanjifamilydefault}{\gtdefault}
\renewcommand{\emph}[1]{{\upshape\bfseries #1}}

\usepackage{amsmath}
\usepackage{amsthm}
\usepackage{tikz}
\usepackage{color}
\usepackage{ascmac}
\usepackage{amsfonts}
\usepackage{mathrsfs}
\usepackage[mathscr]{eucal}
\usepackage{mathtools}
\usepackage{amssymb}
\usepackage{graphicx}
\usepackage{fancybox}
\usepackage{enumerate}
\usepackage{verbatim}
\usepackage{subfigure}
\usepackage{proof}
\usepackage{listings}
\usepackage{otf}
\usepackage[all]{xy}
\usepackage{amscd}
%\usepackage[dvipdfmx]{hyperref}

\usepackage{xcolor}
\definecolor{darkgreen}{rgb}{0,0.45,0} 
\definecolor{darkred}{rgb}{0.75,0,0}
\definecolor{darkblue}{rgb}{0,0,0.6} 
\hypersetup{
    colorlinks=true,
    citecolor=darkgreen,
    linkcolor=darkblue,
    urlcolor=darkblue,
}

\usetikzlibrary{positioning}


\usepackage{latexsym}
\usepackage{wrapfig}
\usepackage{layout}
\usepackage{url}

\usepackage{okumacro}

\usepackage{comment}
%\usepackage{pxjahyper}
% ----------------------------
% commmand
% ----------------------------
% 執筆に便利なコマンド集です. 
% コマンドを追加する場合は下のスペースへ. 

% 集合の記号 (黒板文字)
\newcommand{\NN}{\mathbb{N}}
\newcommand{\ZZ}{\mathbb{Z}}
\newcommand{\QQ}{\mathbb{Q}}
\newcommand{\RR}{\mathbb{R}}
\newcommand{\CC}{\mathbb{C}}
\newcommand{\PP}{\mathbb{P}}
\newcommand{\KK}{\mathbb{K}}


% 集合の記号 (太文字)
\newcommand{\nn}{\mathbf{N}}
\newcommand{\zz}{\mathbf{Z}}
\newcommand{\qq}{\mathbf{Q}}
\newcommand{\rr}{\mathbf{R}}
\newcommand{\cc}{\mathbf{C}}
\newcommand{\pp}{\mathbf{P}}
\newcommand{\kk}{\mathbf{k}}

% 特殊な写像の記号
\newcommand{\ev}{\mathop{\mathrm{ev}}\nolimits} % 値写像
\newcommand{\pr}{\mathop{\mathrm{pr}}\nolimits} % 射影

% スクリプト体にするコマンド
%   例えば　{\mcal C} のように用いる
\newcommand{\mcal}{\mathcal}

% 花文字にするコマンド 
%   例えば　{\h C} のように用いる
\newcommand{\h}{\mathscr}

% ヒルベルト空間などの記号
\newcommand{\F}{\mcal{F}}
\newcommand{\X}{\mcal{X}}
\newcommand{\Y}{\mcal{Y}}
\newcommand{\Hil}{\mcal{H}}
\newcommand{\RKHS}{\Hil_{k}}
\newcommand{\Loss}{\mcal{L}_{D}}
\newcommand{\MLsp}{(\X, \Y, D, \Hil, \Loss)}

% 偏微分作用素の記号
\newcommand{\p}{\partial}

% 角カッコの記号 (内積は下にマクロがあります)
\newcommand{\lan}{\langle}
\newcommand{\ran}{\rangle}



% 圏の記号など
\newcommand{\Set}{{\bf Set}}
\newcommand{\Vect}{{\bf Vect}}
\newcommand{\FDVect}{{\bf FDVect}}
%\newcommand{\Ring}{{\bf Ring}}
\newcommand{\Ab}{{\bf Ab}}
\newcommand{\Mod}{\mathop{\mathrm{Mod}}\nolimits}
\newcommand{\Modf}{\mathop{\mathrm{Mod}^\mathrm{f}}\nolimits}
\newcommand{\CGA}{{\bf CGA}}
\newcommand{\GVect}{{\bf GVect}}
\newcommand{\Lie}{{\bf Lie}}
\newcommand{\dLie}{{\bf Liec}}



% 射の集合など
\newcommand{\Map}{\mathop{\mathrm{Map}}\nolimits} % 写像の集合
\newcommand{\Hom}{\mathop{\mathrm{Hom}}\nolimits} % 射集合
\newcommand{\End}{\mathop{\mathrm{End}}\nolimits} % 自己準同型の集合
\newcommand{\Aut}{\mathop{\mathrm{Aut}}\nolimits} % 自己同型の集合
\newcommand{\Mor}{\mathop{\mathrm{Mor}}\nolimits} % 射集合
\newcommand{\Ker}{\mathop{\mathrm{Ker}}\nolimits} % 核
\newcommand{\Img}{\mathop{\mathrm{Im}}\nolimits} % 像
\newcommand{\Cok}{\mathop{\mathrm{Coker}}\nolimits} % 余核
\newcommand{\Cim}{\mathop{\mathrm{Coim}}\nolimits} % 余像

% その他便利なコマンド
\newcommand{\dip}{\displaystyle} % 本文中で数式モード
\newcommand{\e}{\varepsilon} % イプシロン
\newcommand{\dl}{\delta} % デルタ
\newcommand{\pphi}{\varphi} % ファイ
\newcommand{\ti}{\tilde} % チルダ
\newcommand{\pal}{\parallel} % 平行
\newcommand{\op}{{\rm op}} % 双対を取る記号
\newcommand{\lcm}{\mathop{\mathrm{lcm}}\nolimits} % 最小公倍数の記号
\newcommand{\Probsp}{(\Omega, \F, \P)} 
\newcommand{\argmax}{\mathop{\rm arg~max}\limits}
\newcommand{\argmin}{\mathop{\rm arg~min}\limits}





% ================================
% コマンドを追加する場合のスペース 
%\newcommand{\OO}{\mcal{O}}



\makeatletter
\renewenvironment{proof}[1][\proofname]{\par
  \pushQED{\qed}%
  \normalfont \topsep6\p@\@plus6\p@\relax
  \trivlist
  \item[\hskip\labelsep
%        \itshape
         \bfseries
%    #1\@addpunct{.}]\ignorespaces
    {#1}]\ignorespaces
}{%
  \popQED\endtrivlist\@endpefalse
}
\makeatother

\renewcommand{\proofname}{\textrm{Proof.}}



%\renewcommand\proofname{\bf 証明} % 証明
\numberwithin{equation}{subsection}
\newcommand{\cTop}{\textsf{Top}}
%\newcommand{\cOpen}{\textsf{Open}}
\newcommand{\Op}{\mathop{\textsf{Op}}\nolimits}
\newcommand{\Ob}{\mathop{\textrm{Ob}}\nolimits}
\newcommand{\id}{\mathop{\mathrm{id}}\nolimits}
\newcommand{\pt}{\mathop{\mathrm{pt}}\nolimits}
\newcommand{\res}{\mathop{\rho}\nolimits}
\newcommand{\A}{\mcal{A}}
\newcommand{\B}{\mcal{B}}
\newcommand{\C}{\mcal{C}}
\newcommand{\D}{\mcal{D}}
\newcommand{\E}{\mcal{E}}
\newcommand{\G}{\mcal{G}}
%\newcommand{\H}{\mcal{H}}
\newcommand{\I}{\mcal{I}}
\newcommand{\J}{\mcal{J}}
\newcommand{\OO}{\mcal{O}}
\newcommand{\Ring}{\mathop{\textsf{Ring}}\nolimits}
\newcommand{\cAb}{\mathop{\textsf{Ab}}\nolimits}
%\newcommand{\Ker}{\mathop{\mathrm{Ker}}\nolimits}
\newcommand{\im}{\mathop{\mathrm{Im}}\nolimits}
\newcommand{\Coker}{\mathop{\mathrm{Coker}}\nolimits}
\newcommand{\Coim}{\mathop{\mathrm{Coim}}\nolimits}
\newcommand{\rank}{\mathop{\mathrm{rank}}\nolimits}
\newcommand{\Ht}{\mathop{\mathrm{Ht}}\nolimits}
\newcommand{\supp}{\mathop{\mathrm{supp}}\nolimits}
\newcommand{\colim}{\mathop{\mathrm{colim}}}
\newcommand{\Tor}{\mathop{\mathrm{Tor}}\nolimits}

\newcommand{\cat}{\mathscr{C}}

%筆記体
\newcommand{\cA}{\mcal{A}}
\newcommand{\cB}{\mcal{B}}
\newcommand{\cC}{\mcal{C}}
\newcommand{\cD}{\mcal{D}}
\newcommand{\cE}{\mcal{E}}
\newcommand{\cF}{\mcal{F}}
\newcommand{\cG}{\mcal{G}}
\newcommand{\cH}{\mcal{H}}
\newcommand{\cI}{\mcal{I}}
\newcommand{\cJ}{\mcal{J}}
\newcommand{\cK}{\mcal{K}}
\newcommand{\cL}{\mcal{L}}
\newcommand{\cM}{\mcal{M}}
\newcommand{\cN}{\mcal{N}}
\newcommand{\cO}{\mcal{O}}
\newcommand{\cP}{\mcal{P}}
\newcommand{\cQ}{\mcal{Q}}
\newcommand{\cR}{\mcal{R}}
\newcommand{\cS}{\mcal{S}}
\newcommand{\cT}{\mcal{T}}
\newcommand{\cU}{\mcal{U}}
\newcommand{\cV}{\mcal{V}}
\newcommand{\cW}{\mcal{W}}
\newcommand{\cX}{\mcal{X}}
\newcommand{\cY}{\mcal{Y}}
\newcommand{\cZ}{\mcal{Z}}


\newcommand{\scA}{\mathscr{A}}
\newcommand{\scB}{\mathscr{B}}
\newcommand{\scC}{\mathscr{C}}
\newcommand{\scD}{\mathscr{D}}
\newcommand{\scE}{\mathscr{E}}
\newcommand{\scF}{\mathscr{F}}
\newcommand{\scN}{\mathscr{N}}
\newcommand{\scO}{\mathscr{O}}
\newcommand{\scV}{\mathscr{V}}
\newcommand{\scU}{\mathscr{U}}


\newcommand{\ibA}{\mathop{\text{\textit{\textbf{A}}}}}
\newcommand{\ibB}{\mathop{\text{\textit{\textbf{B}}}}}
\newcommand{\ibC}{\mathop{\text{\textit{\textbf{C}}}}}
\newcommand{\ibD}{\mathop{\text{\textit{\textbf{D}}}}}
\newcommand{\ibE}{\mathop{\text{\textit{\textbf{E}}}}}
\newcommand{\ibF}{\mathop{\text{\textit{\textbf{F}}}}}
\newcommand{\ibG}{\mathop{\text{\textit{\textbf{G}}}}}
\newcommand{\ibH}{\mathop{\text{\textit{\textbf{H}}}}}
\newcommand{\ibI}{\mathop{\text{\textit{\textbf{I}}}}}
\newcommand{\ibJ}{\mathop{\text{\textit{\textbf{J}}}}}
\newcommand{\ibK}{\mathop{\text{\textit{\textbf{K}}}}}
\newcommand{\ibL}{\mathop{\text{\textit{\textbf{L}}}}}
\newcommand{\ibM}{\mathop{\text{\textit{\textbf{M}}}}}
\newcommand{\ibN}{\mathop{\text{\textit{\textbf{N}}}}}
\newcommand{\ibO}{\mathop{\text{\textit{\textbf{O}}}}}
\newcommand{\ibP}{\mathop{\text{\textit{\textbf{P}}}}}
\newcommand{\ibQ}{\mathop{\text{\textit{\textbf{Q}}}}}
\newcommand{\ibR}{\mathop{\text{\textit{\textbf{R}}}}}
\newcommand{\ibS}{\mathop{\text{\textit{\textbf{S}}}}}
\newcommand{\ibT}{\mathop{\text{\textit{\textbf{T}}}}}
\newcommand{\ibU}{\mathop{\text{\textit{\textbf{U}}}}}
\newcommand{\ibV}{\mathop{\text{\textit{\textbf{V}}}}}
\newcommand{\ibW}{\mathop{\text{\textit{\textbf{W}}}}}
\newcommand{\ibX}{\mathop{\text{\textit{\textbf{X}}}}}
\newcommand{\ibY}{\mathop{\text{\textit{\textbf{Y}}}}}
\newcommand{\ibZ}{\mathop{\text{\textit{\textbf{Z}}}}}

\newcommand{\ibx}{\mathop{\text{\textit{\textbf{x}}}}}

%\newcommand{\Comp}{\mathop{\mathrm{C}}\nolimits}
%\newcommand{\Komp}{\mathop{\mathrm{K}}\nolimits}
%\newcommand{\Domp}{\mathop{\mathsf{D}}\nolimits}%複体のホモトピー圏
%\newcommand{\Comp}{\mathrm{C}}
%\newcommand{\Komp}{\mathrm{K}}
%\newcommand{\Domp}{\mathsf{D}}%複体のホモトピー圏

\newcommand{\Comp}{\mathsf{C}}
\newcommand{\Komp}{\mathsf{K}}
\newcommand{\Domp}{\mathsf{D}}
\newcommand{\Kompl}{\mathop{\mathsf{K}^\mathrm{+}}\nolimits}
\newcommand{\Kompu}{\mathop{\mathsf{K}^\mathrm{-}}\nolimits}
\newcommand{\Kompb}{\mathop{\mathsf{K}^\mathrm{b}}\nolimits}
\newcommand{\Dompl}{\mathop{\mathsf{D}^\mathrm{+}}\nolimits}
\newcommand{\Dompu}{\mathop{\mathsf{D}^\mathrm{-}}\nolimits}
\newcommand{\Dompb}{\mathop{\mathsf{D}^\mathrm{b}}\nolimits}
\newcommand{\Dompbf}{\mathop{\mathsf{D}_\mathrm{f}^\mathrm{b}}\nolimits}




\newcommand{\CCat}{\Comp(\cat)}
\newcommand{\KCat}{\Komp(\cat)}
\newcommand{\DCat}{\Domp(\cat)}%圏Cの複体のホモトピー圏
\newcommand{\HOM}{\mathop{\mathscr{H}\hspace{-2pt}om}\nolimits}%内部Hom
\newcommand{\RHOM}{\mathop{\mathrm{R}\hspace{-1.5pt}\HOM}\nolimits}

\newcommand{\muS}{\mathop{\mathrm{SS}}\nolimits}
\newcommand{\RG}{\mathop{\mathrm{R}\hspace{-0pt}\Gamma}\nolimits}
\newcommand{\RHom}{\mathop{\mathrm{R}\hspace{-1.5pt}\Hom}\nolimits}
\newcommand{\Rder}{\mathrm{R}}

\newcommand{\simar}{\mathrel{\overset{\sim}{\rightarrow}}}%同型右矢印
\newcommand{\simarr}{\mathrel{\overset{\sim}{\longrightarrow}}}%同型右矢印
\newcommand{\simra}{\mathrel{\overset{\sim}{\leftarrow}}}%同型左矢印
\newcommand{\simrra}{\mathrel{\overset{\sim}{\longleftarrow}}}%同型左矢印

\newcommand{\hocolim}{{\mathrm{hocolim}}}
\newcommand{\indlim}[1][]{\mathop{\varinjlim}\limits_{#1}}
\newcommand{\sindlim}[1][]{\smash{\mathop{\varinjlim}\limits_{#1}}\,}
\newcommand{\Pro}{\mathrm{Pro}}
\newcommand{\Ind}{\mathrm{Ind}}
\newcommand{\prolim}[1][]{\mathop{\varprojlim}\limits_{#1}}
\newcommand{\sprolim}[1][]{\smash{\mathop{\varprojlim}\limits_{#1}}\,}

\newcommand{\Sh}{\mathrm{Sh}}
\newcommand{\PSh}{\mathrm{PSh}}

\newcommand{\rmD}{\mathrm{D}}

\newcommand{\Lloc}[1][]{\mathord{\mathcal{L}^1_{\mathrm{loc},{#1}}}}
\newcommand{\ori}{\mathord{\mathrm{or}}}
\newcommand{\Db}{\mathord{\mathscr{D}b}}

\newcommand{\codim}{\mathop{\mathrm{codim}}\nolimits}



\newcommand{\gld}{\mathop{\mathrm{gld}}\nolimits}
\newcommand{\wgld}{\mathop{\mathrm{wgld}}\nolimits}


\newcommand{\tens}[1][]{\mathbin{\otimes_{\raise1.5ex\hbox to-.1em{}{#1}}}}
\newcommand{\etens}{\mathbin{\boxtimes}}
\newcommand{\ltens}[1][]{\mathbin{\overset{\mathrm{L}}\tens}_{#1}}
\newcommand{\mtens}[1][]{\mathbin{\overset{\mathrm{\mu}}\tens}_{#1}}
\newcommand{\lltens}[1][]{{\mathop{\tens}\limits^{\mathrm{L}}_{#1}}}
\newcommand{\letens}{\overset{\mathrm{L}}{\etens}}
\newcommand{\detens}{\underline{\etens}}
\newcommand{\ldetens}{\overset{\mathrm{L}}{\underline{\etens}}}
\newcommand{\dtens}[1][]{{\overset{\mathrm{L}}{\underline{\otimes}}}_{#1}}

\newcommand{\blk}{\mathord{\ \cdot\ }}
\newcommand{\mres}[2][]{{\left.{#1}\right\rvert}_{#2}}


%\newcommand{\hocolim}{{\mathrm{hocolim}}}
%\newcommand{\indlim}[1][]{\mathop{\varinjlim}\limits_{#1}}
%\newcommand{\sindlim}[1][]{\smash{\mathop{\varinjlim}\limits_{#1}}\,}
%\newcommand{\Pro}{\mathrm{Pro}}
%\newcommand{\Ind}{\mathrm{Ind}}
%\newcommand{\prolim}[1][]{\mathop{\varprojlim}\limits_{#1}}
%\newcommand{\sprolim}[1][]{\smash{\mathop{\varprojlim}\limits_{#1}}\,}
\newcommand{\proolim}[1][]{\mathop{\text{\rm``{$\varprojlim$}''}}\limits_{#1}}
\newcommand{\sproolim}[1][]{\smash{\mathop{\rm``{\varprojlim}''}\limits_{#1}}}
\newcommand{\inddlim}[1][]{\mathop{\text{\rm``{$\varinjlim$}''}}\limits_{#1}}
\newcommand{\sinddlim}[1][]{\smash{\mathop{\text{\rm``{$\varinjlim$}''}}\limits_{#1}}\,}
\newcommand{\ooplus}{\mathop{\text{\rm``{$\oplus$}''}}\limits}
\newcommand{\bbigsqcup}{\mathop{``\bigsqcup''}\limits}
\newcommand{\bsqcup}{\mathop{``\sqcup''}\limits}
\newcommand{\dsum}[1][]{\mathbin{\oplus_{#1}}}

\newcommand{\Fct}{\mathop{\mathsf{Fct}}\nolimits}
\newcommand{\Red}{\mathop{\mathsf{Red}}\nolimits}
\newcommand{\Cone}{\mathop{\mathsf{Cone}}\nolimits}
\newcommand{\epi}{\mathop{\mathsf{epi}}\nolimits}





%================================================
% 自前の定理環境
%   https://mathlandscape.com/latex-amsthm/
% を参考にした
\newtheoremstyle{mystyle}%   % スタイル名
    {5pt}%                   % 上部スペース
    {5pt}%                   % 下部スペース
    {}%              % 本文フォント
    {}%                  % 1行目のインデント量
    {\bfseries}%                      % 見出しフォント
    {.}%                     % 見出し後の句読点
    {12pt}%                     % 見出し後のスペース
    {\thmname{#1}\thmnumber{ #2}\thmnote{{\hspace{2pt}\normalfont (#3)}}}% % 見出しの書式

\theoremstyle{mystyle}
\newtheorem{AXM}{公理}%[section]
\newtheorem{DFN}[AXM]{定義}
\newtheorem{THM}[AXM]{定理}
\newtheorem*{THM*}{定理}
\newtheorem{PRP}[AXM]{命題}
\newtheorem{LMM}[AXM]{補題}
\newtheorem{CRL}[AXM]{系}
\newtheorem{EG}[AXM]{例}
\newtheorem*{EG*}{例}
\newtheorem{RMK}[AXM]{注意}
\newtheorem{CNV}[AXM]{約束}
\newtheorem{CMT}[AXM]{コメント}
\newtheorem*{CMT*}{コメント}
\newtheorem{NTN}[AXM]{記号}
\newtheorem{CLM}[AXM]{Claim}
\newtheorem{Prop}[AXM]{Proposition}

% 定理環境ここまで
%====================================================

\usepackage{framed}
\definecolor{lightgray}{rgb}{0.75,0.75,0.75}
\renewenvironment{leftbar}{%
  \def\FrameCommand{\textcolor{lightgray}{\vrule width 4pt} \hspace{10pt}}% 
  \MakeFramed {\advance\hsize-\width \FrameRestore}}%
{\endMakeFramed}
\newenvironment{redleftbar}{%
  \def\FrameCommand{\textcolor{lightgray}{\vrule width 1pt} \hspace{10pt}}% 
  \MakeFramed {\advance\hsize-\width \FrameRestore}}%
 {\endMakeFramed}



\renewcommand{\Re}{\mathop{\mathrm{Re}}\nolimits}


% =================================





% ---------------------------
% new definition macro
% ---------------------------
% 便利なマクロ集です

% 内積のマクロ
%   例えば \inner<\pphi | \psi> のように用いる
\def\inner<#1>{\langle #1 \rangle}

% ================================
% マクロを追加する場合のスペース 

%=================================







\def\fracinline#1/#2{\mbox{\raise0.5ex\hbox{\footnotesize$#1$}{\hskip-.1em$/$\hskip-.1em}\raise-0.5ex\hbox{\footnotesize$#2$}}}

\begin{document}




\begin{frame}
    \titlepage
\end{frame}
\begin{comment}
    \begin{frame}\frametitle{書くこと}
    \begin{itemize}
        \item 特に興味を持った定理や理論の背景
        \item それを記述するための記号や概念の導入
        \item 正確な主張の紹介
        \item 証明の基本的なアイデアやアウトラインの紹介
        \item または定理の過程を満たす具体例や定理の応用例の紹介
    \end{itemize}
    \end{frame}
\end{comment}
%\setcounter{section}{3}
%\section*{はじめに}
\begin{frame}
    \frametitle{発表内容}
    \tableofcontents
\end{frame}



\section[背景]{研究の背景}
%\subsection[層]{層理論}
%\subsection{定義}

\begin{frame}
    \frametitle{歴史}
    \begin{description}
        \item[1945前後] Leray，層を定義\cite{Ler45}
        \item[1970前後] 佐藤，柏原・河合らと超局所解析を展開\cite{SKK73}
        \item[1982--1990] 柏原・Schapira, 層の超局所理論を創始（cf. \cite{KS90}）
        \item[2000以降] シンプレクティック幾何への応用が盛んに（cf. \cite{GKS12}等）   
    \end{description}
\end{frame}

\begin{frame}
    \frametitle{層係数コホモロジー：前置き}
    空間\(X\)に対し，そのコホモロジー
    \[
        H^i(X;G)\quad\text{（\(G\)はアーベル群）}
    \]は重要な量．
    
    場所によって係数\(G\)を取り替えると，
    局所系や層を係数とするコホモロジー
    \[
        H^i(X;\mathscr{F})
    \]が出てくる．
    （\(
        \text{局所系}
        \leftrightarrow
        \text{局所定数層}
        \subset\text{層}
    \)）
\end{frame}

\begin{frame}
    \frametitle{層：簡単な例}

    連続関数に対する2つの操作：
    \begin{itemize}
        \item 制限
        \item はりあわせ
    \end{itemize}
    を
    \begin{itemize}
        \item 制限\(\to\)前層
        \item 貼り合わせ\(\to\)層
    \end{itemize}
    として定式化

    \begin{RMK}
        関数環のような層は局所系では拾えない層の例
    \end{RMK}
\end{frame}



\begin{frame}
    \frametitle{層の例}
    \begin{EG}\label{eg:sheaf}
        \(X\)：位相空間，\(Z\subset X\)：閉集合
        
        \(\kk\)：体
        \[
            U\mapsto
            \kk_{X,Z}(U)\coloneqq\left\{
                f\colon Z\cap U\to\kk; 
                \text{\(f\)は局所定数関数}
            \right\}
        \]
    \end{EG}    
    \begin{EG}
        \(E\to X\)：\(C^{\infty}\)ベクトル束
        \[
            C^{\infty}(E)\colon U\mapsto \Gamma(U;E)\coloneqq
            \left\{s\colon U\to E;s\colon C^{\infty}, \pi\circ s=\id_U\right\}
        \]%は層である．
    \end{EG}
    %一般にベクトル束は局所自由層と対応する．
\end{frame}
\begin{frame}
    \frametitle{\(Z\)に台を持つ層のイメージ}
\end{frame}

\begin{frame}
    \frametitle{茎：各点周りに注目}
    \begin{itemize}
        \item 正則関数\(f\in\OO_{X}(U)\)は
        各点\(P\)のまわりで冪級数展開可能
        \item 定義域の異なる関数（環）に対して，
        各点まわりに注目するという操作を
        \begin{center}
            層\(\mathscr{F}\)の\(P\)での茎\(\mathscr{F}_{P}\)
        \end{center}
        として定式化．        
    \end{itemize}
\end{frame}


%\section[層係数コホモロジー]{層係数コホモロジー}


\begin{frame}
    \frametitle{層係数コホモロジー}
    コホモロジーを計算するときには，``良い層''\(\mathscr{I}^i\)たちで分解して
    \[\vcenter{\xymatrix@C=22pt@R=26pt{
        \text{もとの層}\mathscr{F}
        \ar@{~>}[d]^-{\text{分解}}
        \\ 
        \text{層の複体\(\mathscr{I}^\bullet\)}
        \ar@{~>}[d]^-{H^i}
        \\        
        \text{コホモロジー}H^i(X;\mathscr{F})
    }}\]
    のようにする
\end{frame}
\begin{frame}
    \frametitle{層係数コホモロジー}
    もっと一般に，アーベル圏\(\mathscr{A}, \mathscr{B}\)に対して，
    左完全関手\(F\colon \mathscr{A}\to \mathscr{B}\)があるとき，
    \[
        \rr^iF\colon\mathscr{A}\to \mathscr{B}
    \]
    を定めた

    導来圏の理論を用いると見通しがよい
    \begin{NTN}
        層の複体を対象とし，
        コホモロジー間の同型を誘導する射を可逆にした
        圏を\(\Dompb(\kk_X)\)で表す．
    \end{NTN}
\end{frame}
%\section{層の超局所台}
\begin{frame}\frametitle{層の超局所台}
    \(\Dompb(\kk_X)\ni F\mapsto \SS(F)\subset T^{\ast}X\)
    \begin{Definition}[層の超局所台]
        \(p\notin \SS(F)\)となるのは，
        \(\exists U\in I_{p}\)で\(
            \forall x_0\in X,
            \varphi\colon X\to\rr
        \)で\(d\varphi(x_0)\in U\)となるものに対し，\[
            \indlim[x_0\in B]H^{n}(B;F)\simarr
            \indlim[x_0\in B]H^{n}(B\cap \{\varphi<\varphi(x_0)\};F)
        \]
        となるものが存在するとき．
    \end{Definition}
\end{frame}

\begin{frame}
    \frametitle{
        超局所台の例：\(Z\coloneqq\{y\leqq0\}\subset\rr^2\)に
        台を持つ定数層\(\kk_{\rr^2,Z}\)の場合}
\end{frame}

\section[ホモロジカル劣微分]{ホモロジカル劣微分}

\begin{frame}
    \frametitle{古典的な劣微分}
    \(x_0\in \rr^n\)，\(f\colon \rr^n\to\rr\)：凸関数に対し，
    \(f\)の\(x_0\)における劣微分\(\partial f\)を
    \[
        \partial f\coloneqq \{
            \alpha\in\rr^n;
            f(x)>f(x_0)+\langle \alpha, x-x_0\rangle
        \}
    \]
    で定める．
\end{frame}\begin{frame}
    \frametitle{古典的な劣微分のイメージ}
\end{frame}
\begin{frame}
    \frametitle{先行研究}
    \begin{itemize}
        \item 
        \cite{Vic13}は超局所台を用いて，劣微分を下半連続関数に拡張，
        \item 
        \cite{Jub13}はWhitney cone や超局所台を用いて，多様体間の写像の正則性を判定
    \end{itemize}
    以下，\cite{Vic13}による劣微分（ホモロジカル劣微分）の構成について説明
\end{frame}


\begin{frame}
    \frametitle{補助的な定義}

    \begin{itemize}
        \item \(X\)：\(C^\infty\) 多様体
        \item \(J^{1}(X)\coloneqq T^{\ast}X\times \rr\)：1ジェット空間，\((x;\xi,t)\)
        \item \(\tilde{\lambda}=\lambda+dt\)：\(J^{1}(X)\)の接触形式
    \end{itemize}
    
    \[
        \vcenter{\xymatrix@C=26pt@R=26pt{
        \left\{\tau>0\right\}\cap T^{\ast}(X\times\rr)
        \ar[rd]_-{\tilde{\rho}_{t}}
        \ar[rr]^-{\rho_{t}}
        &&
        T^{\ast}X
        \\
        &
        J^{1}(X)
        \ar[ru]_-{r}
      }}
    \]
    \[
        \rho_{t}\left(x,t;\xi,\tau\right)
        \coloneqq
        \left(x,\frac{\xi}{\tau}\right),
        \quad
        \tilde{\rho}_{t}\left(x,t;\xi,\tau\right)
        \coloneqq
        \left(x;\frac{\xi}{\tau},t\right),
        \quad
        r\colon\text{projection}.
    \]

\end{frame}

\begin{frame}
    \frametitle{錐化と簡約}
    \begin{definition}[錐化集合 (conification)]
        \begin{itemize}
            \item \(L\subset T^{\ast}X\)：滑らかなラグランジュ部分多様体
        \end{itemize}
        \[
            \Cone(L)\coloneqq
            \rho_{t}^{-1}(L)\subset 
            T^{\ast}X\times T^{\ast}\rr
            \colon\text{\((n+2)\)次元部分多様体}
        \]
    \end{definition}
    
    \begin{definition}[簡約集合 (reduction)]
        \begin{itemize}
            \item \(A\subset T^{\ast}X\times T^{\ast}\rr\)：錐状部分集合
        \end{itemize}
        \[\Red(A)\coloneqq\rho_{t}(A\cap\left\{\tau>0\right\})\]
    \end{definition}

\end{frame}
\begin{frame}
    \frametitle{錐化と簡約}

    \begin{CLM}
        \(\Red(\Cone(L))=L\subset T^{\ast}X\)
    \end{CLM}
\end{frame}


\begin{frame}
    \frametitle{層の representative}

    \begin{itemize}
        \item \(X\): \(C^\infty\) 多様体
        \item \(\mathbf{k}\): 単位元を持つ大域次元が有限な可換環
    \end{itemize}
    
    \begin{definition}[層の representative]
        \begin{itemize}
            \item \(F\in \Dompb(\kk_{X\times \rr})\)に対し，
            次の集合\(\mathrm{R}(F)\)を\(F\)の representative という．
        \end{itemize}
        \[\mathrm{R}(F)\coloneqq\Red(\SS(F))\subset T^{\ast}X\]
    \end{definition}

\end{frame}

\begin{frame}
    \frametitle{エピグラフ}
    
    \begin{definition}[エピグラフ (epigraph)]
        \begin{itemize}
            \item \(f\colon X\to \rr\)に対し，
            次の集合\(\epi(f)\)を\(f\)のエピグラフという．
        \end{itemize}
        \[
            \epi(f)\coloneqq
            \left\{(x,t)\in X\times \rr;f(x)\leqq t\right\}.
        \]
    \end{definition}
    \(\epi(f)\)を\(Z_f\)ともかく．
    エピグラフに台を持つ層を\(
            F_{f}\coloneqq\kk_{\epi(f)}\)で表す．
    \begin{definition}[ホモロジカル劣微分]
        下半連続関数\(f\colon X\to \rr\)に対し，
        \(f\)の劣微分\(\p f\)とは
        \[\p f\coloneqq \mathrm{R}(F_{f})^{a}\]
    \end{definition}

\end{frame}
\begin{frame}
    \frametitle{構成まとめ}
    劣微分の定義をダイジェスト的に述べると
    \begin{itemize}
        \item \(f\colon X\to \rr\)：下半連続関数
        \item \(\epi(f)\subset X\times \rr\)：\(f\)のエピグラフ
        \item \(F_{f}=\kk_{\epi(f)}\)：エピグラフに台を持つ定数層
        \item \(\SS(F_{f})\subset T^{\ast}X\times T^{\ast}\rr\)：超局所台をとる
        \item \(\p f = -\Red(\SS(F_{f}))\subset T^{\ast}X\)：「射影」（簡約）をとってひねる
    \end{itemize}
\end{frame}

\begin{frame}
    \frametitle{\(C^1\)級関数の劣微分は1点}
    \begin{NTN}
        \begin{itemize}
            \item \(\mres[\partial f]{x}=\partial f\cap T^{\ast}_{x}X\)
            \item \(\partial f(x)=\mres[\partial f]{x}\)
        \end{itemize}
    \end{NTN}
    \begin{EG}
        \(f\in C^{1}(X)\)とすると\[
            \mres[\partial f]{x}=\left\{df(x)\right\}.
        \]
    \end{EG}
\end{frame}

\begin{frame}
    \frametitle{[Jub13]での扱い}
    \(\Gamma_{f}\)：多様体の連続写像\(f\colon M\to N\)のグラフ
    \begin{DFN}
    \begin{itemize}
        \item \(C_f\coloneqq C(\Gamma_f,\Gamma_f)\subset T(M\times N)\)
        \item \(\Lambda_f\coloneqq \muS(\kk_{\Gamma_{f}})\subset T^{\ast}(M\times N)\)
    \end{itemize}
    \end{DFN}
    \begin{THM}[{\cite{Jub13}}]
        \begin{itemize}
            \item \(f\): Lipschitz であることと次は同値
            \item \(C_f\cap (0_M\times TN)\subset 0_{MN}\)
            \item \(\Lambda_f\cap (T^{\ast} M\times 0_{N}^{\ast})\subset 0^{\ast}_{MN}\)
        \end{itemize}
    \end{THM}
\end{frame}

\section{研究の展望}
\begin{frame}
    \frametitle{考えられる問題}
    身近な方面では
    \begin{itemize}
        \item regularlly situated subset (cf. \cite{GS16}) の判定
        \item 超関数の特異性スペクトラムとの関係
        \item H\"order連続の特徴づけ
    \end{itemize}
    応用方面では
    \begin{itemize}
        \item 力学系（Aubry-Mather 理論の拡張）
        \item 最適化問題
        \item 粘性解
    \end{itemize}
\end{frame}


\section{References}

\begin{frame}[allowframebreaks]{References}
    \begin{thebibliography}{90}\beamertemplatetextbibitems
        \bibitem[GKS12]{GKS12} Stephane Guillermou, 
        Masaki Kashiwara, Pierre Schapira, 
        \textit{Sheaf Quantization of Hamiltonian Isotopies 
        and Applications to Nondisplaceability Problems}, 
        Duke Math. J. 161(2), 201--245, (1 February 2012).
        \bibitem[GS16]{GS16} Stephane Guillermou, Pierre Schapira, 
        \textit{Construction of sheaves on the subanalytic site}, 
        Astérisque, Soc. Math. France, 383 1--60 (2016)
        \bibitem[Jub13]{Jub13} Benoit Jubin, 
        \textit{A Microlocal Characterization of Lipschitz Continuity}, 
        \url{https://doi.org/10.4171/PRIMS/54-4-2}.
        \bibitem[KS90]{KS90} Masaki Kashiwara, Pierre Schapira, 
        \textit{Sheaves on Manifolds}, Springer, 1990.
        \bibitem[Ler45]{Ler45} Jean Leray, 
        \textit{Sur la forme des espaces topologiquest sur
        les points fixes de représentations}, J. Math. Pures Appl. 24 (1945), 95--167.
        \bibitem[SKK73]{SKK73} M. Sato, T. Kawai, M. Kashiwara, 
        \textit{Microfunctions and pseudodifrerential equations}, 
        LNM 287, 265--529, Springer-Verlag, 1973.
        \bibitem[Vic13]{Vic13} Nicolas Vichery, 
        \textit{Homological Differential Calculus}, \url{https://doi.org/10.48550/arXiv.1310.4845}.
    \end{thebibliography}
\end{frame}


\begin{comment}
\begin{frame}[noframenumbering]
aaa
\end{frame}
\end{comment}

\end{document}